\subsection{The fresh two-object scope at eighteen arrows}

\begin{lemma}[Proved: a pure categorical rectangular band]
\label{CP2:A-pure-categorical-band}
Let $C$ be a finite category consisting of a categorical
rectangular band with trivial local group and one additional
identity at each object. If $r_i,c_i$ are its positive local
row and column counts, the constant integral left module
has a projective resolution of length at most two. Thus
$H_j(BC;\mathbb Z)=0$ for $j\geq3$.
\end{lemma}
\begin{proof}
Put $R=\mathbb ZC$, choose a minimal idempotent $e$, and
let $P=Re$. At object $i$, $P$ is freely based by the $r_i$
minimal rows. Every band arrow sends every basis element
to its own target row. The augmentation $P\to\mathbf1$
therefore has kernel
\[
 V=\bigoplus_i S_i^{\,r_i-1},
\]
where $S_i$ is the rank-one module supported at $i$ and
annihilated by all band arrows. The representable $R1_i$
surjects onto $S_i$ by taking the coefficient of the extra
identity. Its kernel consists of the band arrows out of $i$.
For each of their $c_i$ source-column labels, the arrows
with that column form a copy of $P$. Hence
\[
 0\longrightarrow P^{c_i}\longrightarrow R1_i
   \longrightarrow S_i\longrightarrow0
\]
is exact over $\mathbb Z$. Splicing these sequences into
$0\to V\to P\to\mathbf1\to0$ gives
\[
 0\longrightarrow\bigoplus_i P^{c_i(r_i-1)}
 \longrightarrow\bigoplus_i(R1_i)^{r_i-1}
 \longrightarrow P\longrightarrow\mathbf1\longrightarrow0.
\]
All displayed terms preceding $\mathbf1$ are projective.
Tensoring with the constant right module proves the claimed
homology vanishing.
\end{proof}

\begin{proposition}[Proved: wholly strong two-object categories through eighteen]
\label{CP2:A-strong-two-eighteen}
A strongly connected two-object category with at most
eighteen arrows cannot have type $\CP^2$.
\end{proposition}
\begin{proof}
An object retract can be deleted, leaving a monoid with
fewer than twenty-three elements, excluded by
Theorem~\ref{A:monoid-minimum23-classification}. Otherwise
both identities are outside the categorical minimal ideal.
Write $r,c$ for its total row and column counts and $G$ for
its local group. If $r=2$ or $c=2$, the local counts on that
side are all one and Lemma~\ref{CP2:A-high-Rees} gives
finite-group homotopy type. In the other cases
$|G|rc+2\leq18$ and $r,c\geq3$, so $G=1$.
If one total is three, its two positive local counts are
one and two, and Lemma~\ref{CP2:A-baseline-binary-representable}
annihilates positive rational products. The sole remaining
possibility has $r=c=4$ and exactly eighteen arrows. Its
sixteen minimal-band arrows and two identities exhaust the
category. Lemma~\ref{CP2:A-pure-categorical-band} gives
$H_4=0$, completing the reduction.
\end{proof}

\begin{proposition}[Proved: rational three-state cones through order twelve]
\label{CP2:A-twelve-three-state-rational}
An ordinary act cone over a monoid of order at most twelve,
on at most three states, cannot have the rational homology
and cohomology ring of $\CP^2$.
\end{proposition}
\begin{proof}
Orders at most eleven follow from
Theorem~\ref{CP2:A-bounded-eleven-acts}. At order twelve,
the relative and two-row act obstructions force at least
three minimal rows and two minimal columns. The local group
is trivial, since a nontrivial one would already require
$3\cdot2\cdot2+1=13$ elements. The possible band shapes and
outside counts are
\[
 (3,2,6),\quad(4,2,4),\quad(5,2,2),\quad(3,3,3).
\]
The first case is excluded by
Proposition~\ref{CP2:A-three-row-six-outside-acts}; the next
two by Lemma~\ref{CP2:A-two-column-four-outside-bounded}.

In the last case there must be exactly three states, since
the binary act obstruction is unconditional. The vanishing
of first homology and a putative nonzero square imply
constant band action by Lemma~\ref{CP2:A-four-state-band}.
Actual constant-action compression gives a monoid with the
same outside algebra $A$, of dimension three, and two
augmentation modules of dimension two. The compressed
monoid has the same rational homology as the given cone,
so its Tor groups are bounded and its second Tor is
one-dimensional. If $A$ is local,
Lemma~\ref{CP2:A-local-three-small-module} excludes this.
If $A$ is nonlocal, Lemma~\ref{A:nonlocal-four} excludes it.
All comparisons just used are actual maps or homotopy
equivalences; no boundedness assertion is inferred from
an isolated degree-two comparison.
\end{proof}

\begin{proposition}[Proved: a corner of order at most two through eighteen]
\label{CP2:A-two-object-small-corner-eighteen}
A two-object collage with at most eighteen arrows and a
corner of order at most two cannot have type $\CP^2$.
\end{proposition}
\begin{proof}
Take the small corner as the left corner $N$. A trivial
corner is an ordinary act cone and is excluded by
Corollary~\ref{CP2:A-ordinary-act-eighteen}. Write $m=|M|$
and $s=|S|$ when $|N|=2$, so $m+s\leq16$.
If $m\leq11$, the trivial right rational $N$ module is
projective, and actual exact-orbit comparison gives a
rational ring equivalence to an ordinary $M$ act cone.
Its homology is bounded under the proposed target, so
Theorem~\ref{CP2:A-bounded-eleven-acts} applies. If
$m\geq13$, there are at most three connecting arrows,
excluded by the universal one- and two-cross obstructions
and Theorem~\ref{CP2:A-three-cross-reduction}.

Only $(m,|N|,s)=(12,2,4)$ remains. If $N=\{1,e\}$ with
$e^2=e$, its actor is a biset retraction onto $eS$. The
cut-retract theorem preserves an actual $\CP^2$ type on
the image. The $N$ action there is trivial and $BN$ is
contractible, so collapsing that corner gives an ordinary
$M$ act cone on at most four states by an actual homotopy
equivalence, a contradiction.

If $N=C_2$ and its action on $S$ is nontrivial, its orbit
set has at most three states. Exact rational orbit comparison
gives a rational ring and homology equivalence to the
ordinary $M$ cone on that set, contradicted by
Proposition~\ref{CP2:A-twelve-three-state-rational}. If the
$C_2$ action is trivial, collapsing $C_2$ gives an actual
categorical retract onto the four-state ordinary cone and
a rational homology equivalence. A retract of an integral
$\CP^2$ model has split integral homology and trivial
fundamental group. This rational equivalence forces its
second and fourth split summands to be $\mathbb Z$, and
pullback transfers the primitive square. Integral recognition
therefore makes the retract an actual $\CP^2$ model, again
contradicting the ordinary-act theorem. The rational and
actual comparisons have thus been kept separate in every
case.
\end{proof}

\begin{proposition}[Proved: the complete reduced eighteen-arrow scope]
\label{CP2:A-two-object-eighteen-scope}
To exclude every category with at most two objects and at
most eighteen arrows, it suffices to exclude the following
twenty-five positive-source collage profiles:
\[
\begin{gathered}
 (5,4,9),\ (5,5,8),\ (5,6,7),\ (5,7,6),\ (5,8,5),\ (5,9,4),\\
 (6,4,8),\ (6,5,7),\ (6,6,6),\ (6,7,5),\ (6,8,4),\\
 (7,3,8),\ (7,4,7),\ (7,5,6),\ (7,6,5),\ (7,7,4),\\
 (8,3,7),\ (8,4,6),\ (8,5,5),\ (8,6,4),\\
 (9,3,6),\ (9,4,5),\ (9,5,4),\\
 (10,3,5),\ (10,4,4).
\end{gathered}
\]
The first corner must have nonzero rational second homology.
This is a complete scope reduction, not a completion claim
for these finite profiles.
\end{proposition}
\begin{proof}
All budgets through seventeen are excluded by
Theorem~\ref{CP2:A-two-object-seventeen}; one-object models
are excluded by the monoid minimum. A remaining two-object
model is connected. Its wholly strong support is excluded
by Proposition~\ref{CP2:A-strong-two-eighteen}; otherwise it
is a collage. If both corner second cohomology groups vanish,
Lemma~\ref{CP2:A-collage-relative-cup} kills its second square.
By opposition choose a positive corner as the source and
write its order $m\geq5$, opposite order $n\geq1$ and
connecting-set size $s\geq1$. Then $m+n+s=18$.
There are exactly seventy-eight such ordered triples.

The preceding small-corner result removes $n\leq2$.
The universal one-, two- and three-cross results remove
$s\leq3$. The remaining inequalities $m\geq5,n\geq3,s\geq4$
leave twenty-eight triples. The cases $(5,3,10)$ and
$(6,3,9)$ are removed by
Theorem~\ref{CP2:A-baseline-right-zero-three-six} when the
opposite corner is the nonprojective-right three-element
monoid $R_3$, and by exact orbit comparison followed by
bounded-eleven for every other three-element corner.
The case $(11,3,4)$ is excluded by
Proposition~\ref{CP2:A-right-zero-four-states} for $R_3$,
and by the same projective endpoint comparison otherwise.
These are exactly the three removed triples, leaving the
displayed twenty-five. The independently executable integer
inventory in \path{monoid/code/two_object18_inventory.py}
records all seventy-eight budgets and their precise status.
\end{proof}

\begin{remark*}[Finite inputs still required at this frontier]
The frozen seventeen-arrow streams do not certify any of
the twenty-five new finite profiles. Positive source catalogs
through order nine are complete; source order ten is newly
required only for $(10,3,5)$ and $(10,4,4)$.

At order ten, the band possibilities are $2\times2$ with
six outside elements, $2\times3$ and $3\times2$ with four,
$2\times4$ and $4\times2$ with two, and $3\times3$ with
one. A nontrivial minimal group can only be $C_2$ with
shape $2\times2$ and two outside elements. Every table of
these shapes must be included before a justified homology
filter; the order-nine positive catalog does not cover them.

The opposite order-nine corner occurs only in $(5,9,4)$.
The acyclic-opposite P5 four-state lemma and exact
projective-right orbit comparison remove respectively
one-row and one-column minimal ideals. The remaining
two-sided order-nine catalog contains $21918$ band tables
and two $C_2$ Rees tables, all $21920$ of which are needed,
including tables of zero second rational homology.
Opposite order eight occurs in $(5,8,5)$ and $(6,8,4)$.
For $(5,8,5)$, Corollary~\ref{CP2:A-P5-eight-five-scope}
removes the one-row cases as well, so the complete $1205$
two-sided-band tables suffice. For $(6,8,4)$, those tables
alone do not yet cover the opposite-corner scope: the
omitted one-row cases require a further proof or their
complete catalogs.
No complete eighteen-arrow exclusion is asserted here.
\end{remark*}
